\section{C. The local Kummer cover and its global obstruction}

We distinguish a finite holomorphic map with smooth source from a finite
map with normal, possibly singular, source. A branched cover of compact
complex manifolds has the former meaning in this section. Write
$V_{\rm tor}=\Z\langle\widehat w,\widehat\delta\rangle$.

\paragraph{Theorem C.1.1 (Proved; checked local Kummer theorem).}
Let $N\geq1$. On the cusp toric cover with period correction $C(q)$,
put $q=v^N$ and $C_N(v)=C(v^N)/N$. The map
\[
 P_N(\xi_1,\xi_2,v)=(\xi_1^N,\xi_2^N,v^N)
\]
descends to a finite flat holomorphic map
$\kappa_N:N_0^{\langle N\rangle}\to N_0$ of degree $N^3$ between
smooth threefolds. The source central fibre is reduced and irreducible,
with normalization $\mathrm{dP}_6$. Its ramification divisor is
$(N-1)W_N$, and $\kappa_N^*W=NW_N$.

The base change $N_0[N]=N_0\times_{\Delta_q}\Delta_v$ is already normal.
The induced finite map $Q_N:N_0^{\langle N\rangle}\to N_0[N]$ has
degree $N^2$. For $N>1$ it is quasi-etale, and its nonflat and non-etale
loci are precisely the inverse images of the double and triple strata.
It is not a normalization morphism. On a regular fibre its image lattice is
\[
 L_N=\Z\langle\widehat\gamma,\widehat u,
                   N\widehat w,N\widehat\delta\rangle.
 \tag{C.1.1}
\]

\begin{proof}
On the target torus the period action is
$\Psi_\lambda(x,q)=(e^{2\pi iC(q)\lambda}q^{B_0\lambda}x,q)$,
where $B_0=\left(\begin{smallmatrix}0&1\\-1&0\end{smallmatrix}\right)$.
The source uses the same fan and the action
\[
 \Psi^N_\lambda(\xi,v)=
 \big(e^{2\pi iC(v^N)\lambda/N}v^{B_0\lambda}\xi,v\big).
\]
The integral shear preserves the $A_2$ triangulation and the
holomorphic unit factor acts in every toric chart. Additivity in $\lambda$
proves the action law. Raising the two torus coordinates to the $N$th
power proves $P_N\Psi^N_\lambda=\Psi_\lambda P_N$ exactly.
There is no root-selection ambiguity: the logarithm being divided by $N$
is the specified holomorphic function $2\pi iC(v^N)\lambda$.

In a maximal smooth toric chart the map is
$(w_0,w_1,w_2)\mapsto(w_0^N,w_1^N,w_2^N)$. Its convergent power-series
ring is free over the target ring on the $N^3$ monomials whose three
exponents lie in $[0,N-1]$. Thus $P_N$ is finite flat. These charts cover
the fan, and their inverse images are the corresponding charts because
scalar multiplication preserves every cone. Freeness and proper
discontinuity of the source lattice action now follow directly from the
target action: a source fixed point would map to a target fixed point;
if two source compact sets meet under a translate, their target compact
images meet under the same translate. Only finitely many such translates
exist. This transfers the established cusp action property without a new
estimate on $C_N$.

Over an evenly covered target open set choose one lift to the toric
cover. Every source lattice orbit over it has a unique representative
whose image is in that lift. Hence the descended morphism is locally the
same finite flat map, of degree $N^3$. It is proper. Since $f_0$ is proper,
the source map to $\Delta_v$ is proper as well: the inverse image of a
compact $K\subset\Delta_v$ is closed in
$\kappa_N^{-1}(f_0^{-1}(\{v^N:v\in K\}))$, a compact set.
The free action on a smooth toric space gives a smooth source. The fan and
its period action on vertices have not changed, so its central fibre has
one component and the same hexagonal normalization. In this chart
$q=z_0z_1z_2$, $v=w_0w_1w_2$, and the Jacobian is
$N^3(w_0w_1w_2)^{N-1}$, proving the asserted divisors.

At a triple point the relative factor has the exact ring map
\[
 \frac{\C\{z_0,z_1,z_2,v\}}{(z_0z_1z_2-v^N)}
 \longrightarrow\C\{w_0,w_1,w_2\},\qquad
 z_i\longmapsto w_i^N,\quad v\longmapsto w_0w_1w_2.
 \tag{C.1.2}
\]
The source ring on the left is the invariant ring for
$G=\{(\epsilon_0,\epsilon_1,\epsilon_2)\in\mu_N^3:
\epsilon_0\epsilon_1\epsilon_2=1\}$ on the right. Indeed an invariant
monomial has three exponents congruent modulo $N$; factor out their
common residue as a power of $w_0w_1w_2$. The same expansion for convergent
series proves the analytic assertion. A finite-group invariant ring of a
normal domain is normal: an element of its fraction field integral over
the invariants is integral over the original ring, belongs to that ring
by normality, and is invariant. This proves normality here and on the
face charts, and the degree is $|G|=N^2$.

The stabilizer is trivial when at most one $w_i$ vanishes, has order $N$
when exactly two vanish, and is $G$ at the origin. The origin fibre has
length
\[
 N^3-(N-1)^3=3N^2-3N+1>N^2\quad(N>1).
\]
At a double point the two-variable factor has generic rank $N$ and
closed fibre length $2N-1>N$. These counts follow by excluding monomials
divisible by the product of the variables from the box of exponents
$0,\ldots,N-1$. They prove failure of flatness exactly where claimed;
elsewhere the free finite action gives etaleness and flatness.

Finally, in additive coordinates multiplication by $N$ is the map
\[
 \C^2/(N\Z^2+Z(Nr)\Z^2)\longrightarrow
 \C^2/(\Z^2+Z(Nr)\Z^2),\qquad v=e^{2\pi ir}.
\]
It proves (C.1.1), degree $N^2$, and the fibre deck group
$(\Z/N)^2$. With $S_N=\operatorname{diag}(1,1,N,N)$ and
$M_0=I+R$, $R^2=0$, direct multiplication gives
$S_N^{-1}M_0^NS_N=M_0$. This also checks the marking in the local
workbench theorem, rather than identifying covers from their degrees
alone.
\end{proof}

\paragraph{Lemma C.1.2 (Proved; smoothness of a cover of the entire fan).}
Let $\Sigma$ be the complete periodic $A_2$ triangulation at height one
used for the cusp, in $N'=\Z^2\oplus\Z$. Suppose a finite connected
holomorphic cover of its toric neighbourhood is etale over
$(\C^*)^2\times\Delta^*$ and has smooth total space. The cocharacter
sublattice defining its restriction to this torus is $nN'$ for some
$n\geq1$. In particular its intersection with the two-dimensional
vanishing lattice is $n\Z^2$.

\begin{proof}
A connected finite cover of $(\C^*)^2\times\Delta^*$ is classified by
a finite-index sublattice $K\subset N'$, since this space has fundamental
group $\Z^3$. The unique normal extension across the toric boundary is
the toric map with the same real cones and lattice $K$. This can be seen
on each affine chart by taking the integral closure of its monomial ring
in the finite torus field extension; the resulting monomials have
exponents in the dual of $K$. The convergent analytic version has the
same local finite quotient charts. Uniqueness of normal extension glues
these charts and applies to the given finite cover.

For a target ray $r_v=(v,1)$ let $a_vr_v$ be its primitive generator in
$K$. A smooth source chart over a triangle with vertices $v_0,v_1,v_2$
requires $a_0r_0,a_1r_1,a_2r_2$ to be a basis of $K$. This is the usual
smoothness criterion, also obtained directly by identifying that chart's
semigroup with the free monoid on its three dual basis vectors.
In the adjacent triangle across the edge $(r_1,r_2)$ the other vertex is
$r_3=r_1+r_2-r_0$. Both target triples are integral bases. Equality of
their indices gives
\[
 [N':K]=a_0a_1a_2=a_3a_1a_2,
 \qquad a_3=a_0.
\]
Expressing $a_3r_3$ in the first source basis gives coefficients
$-1,a_0/a_1,a_0/a_2$. These are integers. Applying the same argument
across the other two edges of the first triangle gives all the reverse
divisibilities. Thus $a_0=a_1=a_2=n$, and their three vectors generate
$K=nN'$. This already proves the result; connectivity of the triangles
also shows directly that every ray has ramification index $n$.

The assertion is about the entire cusp fan. A cover can be smooth in
one triangle and singular in its neighbour, so testing one triple point
cannot replace this argument.
\end{proof}

\paragraph{Theorem C.1.3 (Proved; no smooth global Kummer extension).}
For $N>1$, the local map $\kappa_N$ in Theorem C.1.1 cannot be the
restriction, even to one component of the inverse image of a cusp
neighbourhood, of a finite holomorphic map $Y\to X$ from a connected
smooth compact complex manifold. Additional cusp sheets and additional
vertical branch fibres do not remove the obstruction. For $N=1$ the
displayed local map is the identity and has the identity global extension.

\begin{proof}
Suppose such a finite map $F:Y\to X$ exists. Form the Stein factorization
of the proper map $fF$, with connected fibres $g:Y\to C'$ and finite map
$h:C'\to B$. Near the specified cusp point, the fibres of $g$ are tori
and their maps to the corresponding fibres of $f$ are etale isogenies,
with image lattice $L_N$ of (C.1.1). No component of the ramification
divisor of $F$ can dominate $C'$: a proper dominating component would
meet every nearby regular fibre in this open set where $F$ is etale.
Thus all ramification is vertical. Delete its finitely many base values,
the branch values of $h$, and $0,1,\infty$. Over the remaining connected
base the map is a cover of torus bundles. On every sheet the image
of a fibre fundamental group is a finite-index sublattice of $\Lambda$,
and continuation along a loop with monodromy $A$ carries it to $AL_N$.
Translations do not change this subgroup, since the fibre group is abelian.

Follow a based meridian of the second finite fibre, choosing it to avoid
the additional deleted points, and then return to the same small
punctured cusp disc. The resulting sheet has image lattice $A_2L_N$.
It lies in some component of $F^{-1}(N_0)$; it need not be the original
component. Pull that component back to the unramified toric covering
$Y_\Sigma\to N_0$. This is still a finite cover with smooth total space,
etale off the central fibre. Each connected component therefore satisfies
Lemma C.1.2. Consequently its fibre image lattice must meet
$V_{\rm tor}$ in a scalar sublattice $nV_{\rm tor}$.

The actual second monodromy gives
\[
 A_2\widehat u=\widehat w,\qquad
 A_2\widehat w=-\widehat u+\widehat\delta,\qquad
 A_2\widehat\delta=\widehat\delta.
\]
The first coordinate of $A_2\widehat\gamma$ is one. Taking the
intersection of $A_2L_N$ with $V_{\rm tor}$ therefore gives exactly
\[
 A_2L_N\cap V_{\rm tor}
   =\Z\widehat w\oplus N\Z\widehat\delta.
 \tag{C.1.3}
\]
It cannot equal $nV_{\rm tor}$ when $N>1$: the primitive first generator
forces $n=1$ but the second generator then excludes
$\widehat\delta$. This contradiction proves the theorem. Reversing the
based meridian replaces $A_2$ by its inverse and gives the same
non-scalar intersection, so the obstruction does not depend on that
orientation convention.
\end{proof}

\paragraph{Proposition C.1.4 (Proved; a separate degree obstruction).}
Even allowing a normal singular source, there is no finite global
extension whose whole inverse image of $N_0$ is the single prescribed
map $\kappa_N$, for $N>1$. The monodromy orbit of $L_N$ has at least
$N^2$ elements, whereas such an extension would have a Stein base cover
of degree $N$.

\begin{proof}
Put $D=A_2^2A_1A_2^2A_1^{-1}$. Exact multiplication gives
\[
 D=\begin{pmatrix}1&0&0&0\\0&1&0&0\\6&0&1&0\\-6&1&0&1\end{pmatrix}.
\]
The $N^2$ lattices $D^aM_0^bL_N$, $0\leq a,b<N$, are distinct.
Indeed their intersections with $\ker\gamma$, reduced modulo $N\Lambda$,
are generated respectively by
$\widehat u+b\widehat w+a\widehat\delta$; the coefficient of
$\widehat u$ is one, so equality of these cyclic submodules forces
equality of both residues $a,b$. A base covering supporting one fibre
lattice must have at least as many sheets as its monodromy orbit.
The unique prescribed cusp has local connected-fibre parameter $v$
with $q=v^N$ and is the entire preimage, so the base degree would be
$N$. This contradicts $N^2>N$.

The exact orbit enumeration for $2\leq N\leq8$ is supplied in
\texttt{code/kummer\_orbits.py}; it gives respectively
$12,36,96,150,432,392,768$ elements. These finite computations check the
formulas in small cases. The preceding argument proves the required
lower bound for every $N$ and does not infer it from a finite search.
\end{proof}

\paragraph{Theorem C.1.5 (Proved; normal singular globalizations).}
If a normal compact complex source is allowed, the prescribed local
Kummer cover does extend as one component above the cusp, for every $N$.
There is a connected finite flat such cover, branched only over
$S_1\cup S_2\cup W$, of degree
\[
 N^2|G_N|,\qquad
 G_N=\langle(A_1,-1),(A_2,0)\rangle
 \subset\operatorname{GL}_4(\Z/N)\times\Z/N.
 \tag{C.1.4}
\]
Here matrices and the second coordinates are reduced modulo $N$.
The component in question has degree $N^3$ over $N_0$ and is precisely
$\kappa_N$. For $N>1$ other cusp components include singularities;
this is a cover by a normal space, not by a complex manifold.

\begin{proof}
Remove the three special fibres and write the result as $X^*\to B^*$.
The base retracts to a graph. We choose a continuous section that agrees
with the zero section of the marked torus family on a cusp collar.
Such a choice exists: the only obstruction is in
$H^2(B^*,\text{cusp collar};\Lambda)$, which is zero since the based
cusp loop is a free factor in the fundamental group of the punctured
sphere; equivalently, extend the given section along a relative graph
retraction. The resulting group presentation is
\[
 \pi_1(X^*)=\Lambda\rtimes F(a_1,a_2),\qquad
 a_0=(a_1a_2)^{-1},\quad a_j\lambda a_j^{-1}=A_j\lambda.
\]
The restriction to the cusp collar is the specified splitting with
monodromy $M_0$. A topological finite covering of a complex manifold
has a unique lifted holomorphic atlas, so it is enough first to
construct its subgroup.

Map the free group to $G_N$ by the two generators in (C.1.4) and
map the fibre lattice to $(\Z/N)^4$. This gives an epimorphism
\[
 \pi_1(X^*)\longrightarrow
 K_N=(\Z/N)^4\rtimes G_N,
\]
where only the matrix coordinate acts on the first factor. Let
$U_N\subset(\Z/N)^4$ be the subgroup spanned by $g,u$, of order
$N^2$, and take the connected cover corresponding to the inverse
image of $U_N$. Its degree is $|K_N|/|U_N|=N^2|G_N|$.
An element $\lambda a_0^k$ of the cusp group lies in this subgroup
exactly when $k\equiv0\pmod N$ and
$\lambda\in L_N$. Indeed $a_0$ has second coordinate $1$, and
$M_0^N\equiv I\pmod N$ because $(M_0-I)^2=0$.
Consequently its cusp subgroup is exactly
$L_N\rtimes\langle a_0^N\rangle$. The power map of C.1.1 has
this subgroup, using the zero-section lift of the base loop.
The selected connected collar cover is therefore holomorphically
isomorphic to the punctured part of $\kappa_N$.

We construct the extension, including its normality and flatness.
At any point of the omitted divisor choose a polydisc with divisor
$z_1\cdots z_r=0$, $1\leq r\leq3$. Each connected component of
the restricted finite cover is classified by a finite-index subgroup
$K\subset\Z^r$. Write $M'=\operatorname{Hom}(K,\Z)$, containing
$\Z^r$ by restriction. Its completion is the normal toric analytic
space with ring generated by convergent monomials with exponents
\[
 M'\cap\R_{\geq0}^r
\]
over the remaining polydisc coordinates. This semigroup is saturated,
so the ring is normal. For each coset in $M'/\Z^r$ choose its unique
representative whose coordinates belong to $[0,1)$. The corresponding
monomials form a free basis over the original polydisc ring: subtracting
the representative leaves nonnegative integer coordinates, and distinct
cosets cannot cancel. Thus the extension is finite flat. It is also
the integral closure in the finite covering's local sheaf of meromorphic
functions; we do not use the field of global meromorphic functions on
the compact nonalgebraic space.
Integral closure is unique, so these local extensions glue on all
overlaps, including the self-identifications of $W$.

The resulting space is normal and connected; finiteness over compact
$X$ makes it compact. The same uniqueness identifies the selected
cusp completion with the already normal smooth source in C.1.1.
Theorem C.1.3 forces at least one other cusp component to be singular
when $N>1$. This proves the distinction between the two source
categories and the asserted degree. Exact enumeration gives
$|G_2|=48$, $|G_3|=216$, $|G_4|=3072$; the accompanying program
checks these finite examples, while the construction proves the
statement for arbitrary $N$.
\end{proof}

The local assertions in the supplied workbench are valid, with their
two different degrees. Non-invariance of its fibre sublattice alone
does not prohibit a further base cover. The entire-fan argument excludes
smooth global sources; the finite quotient in C.1.5 gives the different
answer for normal singular sources with additional cusp sheets.
